\chapter{Usage Guide}

\section{Getting Started}

\subsection{Prerequisites}

\textbf{Required Software:}
\begin{itemize}
    \item Microsoft Excel (2016 or later)
    \item Python 3.10+
    \item Git (for version control)
\end{itemize}

\textbf{Required Python Packages:}
\begin{lstlisting}[style=pythonstyle]
# Install via pip
pip install pandas numpy matplotlib xlwings statsmodels
\end{lstlisting}

\subsection{Project Setup}

\begin{enumerate}
    \item Clone or download project to local directory
    \item Update paths in \texttt{Control\_Panel\_New.xlsm} VBA:
    \begin{lstlisting}[style=vbastyle]
Const PYTHON_PATH As String = "C:\Path\To\Python\python.exe"
Const PROJECT_PATH As String = "C:\Path\To\Project\"
    \end{lstlisting}
    
    \item Verify Python environment:
    \begin{lstlisting}
python --version
python -c "import pandas, xlwings, statsmodels"
    \end{lstlisting}
    
    \item Test Control Panel: Open \texttt{Control\_Panel\_New.xlsm}, enable macros
\end{enumerate}

\section{Workflow 1: Import and Map Payroll Data}

\subsection{Step-by-Step Process}

\textbf{Step 1: Prepare Source Workbook}

\begin{enumerate}
    \item Open your payroll Excel file (e.g., \texttt{Pay\_emploier\_sept25.xlsm})
    \item Ensure sheets are named by month (e.g., "January", "February", etc.)
    \item Verify structure:
    \begin{itemize}
        \item Row 7: Employee IDs in columns C onwards
        \item Column A: Salary item codes (optional)
        \item Column B: Salary item labels (required)
        \item Columns C+: Employee data
    \end{itemize}
\end{enumerate}

\textbf{Example Source Sheet:}

\begin{table}[H]
\centering
\footnotesize
\begin{tabular}{|l|l|r|r|r|}
\hline
\textbf{Code} & \textbf{Label} & \textbf{EMP-001} & \textbf{EMP-002} & \textbf{EMP-003} \\
\hline
S001 & SALAIRE\_BASE & 15000 & 12000 & 18000 \\
S002 & PRIME\_TRANSPORT & 500 & 500 & 800 \\
S003 & INDEMNITE & 200 & 150 & 300 \\
... & ... & ... & ... & ... \\
\hline
\end{tabular}
\caption{Source Sheet Structure (Transposed Format)}
\end{table}

\textbf{Step 2: Launch Import Dialog}

\begin{enumerate}
    \item Open \texttt{Control\_Panel\_New.xlsm}
    \item Click \textbf{Button 1: Import \& Map Data}
    \item Wait for GUI to appear (may take 10-15 seconds)
\end{enumerate}

\textbf{Step 3: Select Source Workbook}

\begin{enumerate}
    \item Click "Browse..." button
    \item Navigate to your payroll workbook
    \item Select file and click "Open"
    \item Path will appear in dialog
\end{enumerate}

\textbf{Step 4: Define Mappings}

\begin{enumerate}
    \item GUI shows destination column names (left side)
    \item For each destination column:
    \begin{itemize}
        \item Click "Define Formula" button
        \item Enter formula in popup (e.g., \texttt{+SALAIRE\_BASE label})
        \item Click "Save"
    \end{itemize}
    
    \item Formula syntax examples:
    \begin{itemize}
        \item Single source: \texttt{+SALAIRE\_BASE label}
        \item Sum multiple: \texttt{+BASE label +PRIME label}
        \item Subtraction: \texttt{+BRUT label -COT label}
        \item Multiplication: \texttt{+HEURES label *TAUX label}
    \end{itemize}
\end{enumerate}

\textbf{Step 5: Execute Mapping}

\begin{enumerate}
    \item After defining all formulas, click \textbf{"Execute Mapping"}
    \item Progress bar appears (may take 30-60 seconds)
    \item Success message: "Mapping complete!"
    \item Close dialog
\end{enumerate}

\textbf{Step 6: Verify Output}

\begin{enumerate}
    \item Click \textbf{Button 5: Open Mapped Workbook}
    \item \texttt{data\_sorted.xlsx} opens automatically
    \item Verify each month sheet has correct structure
    \item Check calculations are correct
\end{enumerate}

\subsection{Troubleshooting}

\begin{table}[H]
\centering
\footnotesize
\begin{tabularx}{\textwidth}{lX}
\toprule
\textbf{Error} & \textbf{Solution} \\
\midrule
"No CustomMap sheet" & Source workbook needs a CustomMap sheet (will be created automatically on first save) \\
"Key not found" & Check label spelling in formula matches source sheet exactly \\
"No employees detected" & Verify Row 7 has employee IDs in columns C onwards \\
GUI doesn't appear & Check Python path in VBA, run manually: \texttt{python scripts/run\_excel.py} \\
Slow performance & Close other Excel files, use "Modified Only" option \\
\bottomrule
\end{tabularx}
\caption{Common Import/Map Errors}
\end{table}

\section{Workflow 2: Visualize Payroll Analytics}

\subsection{Step-by-Step Process}

\textbf{Step 1: Ensure Data Exists}

\begin{enumerate}
    \item Verify \texttt{outputs/payroll\_long.csv} exists
    \item If not, run Workflow 1 first or manually convert data
    \item Check file has required columns: \texttt{month}, \texttt{employee\_id}, \texttt{salaire\_brut}, \texttt{salaire\_net}
\end{enumerate}

\textbf{Step 2: Launch Dashboard}

\begin{enumerate}
    \item Open \texttt{Control\_Panel\_New.xlsm}
    \item Click \textbf{Button 4: Visualize Data}
    \item Dashboard GUI appears (10-15 seconds)
\end{enumerate}

\textbf{Step 3: Choose Visualization}

\begin{enumerate}
    \item Dashboard shows 10 buttons:
    \begin{enumerate}
        \item Total Cost Trend
        \item Cost Forecast (AR)
        \item Cost Per Employee
        \item Cost Breakdown
        \item Cost Distribution
        \item Top Employees
        \item Gross vs Net Comparison
        \item Headcount Trend
        \item Burn Rate Dashboard
        \item Year-over-Year Comparison
    \end{enumerate}
    
    \item Click any button to generate that plot
    \item Plot appears in separate matplotlib window
    \item Close plot window to return to dashboard
\end{enumerate}

\textbf{Step 4: Interactive Plots}

For \textbf{Plot 7 (Gross vs Net)}:
\begin{enumerate}
    \item Click "7. Gross vs Net Comparison"
    \item Selection dialog appears with employee list
    \item Scroll and select an employee
    \item Click "Plot" button
    \item Plot shows gross vs net for that employee over time
\end{enumerate}

\subsection{Interpretation Guide}

\textbf{Plot 1: Total Cost Trend}
\begin{itemize}
    \item \textbf{Look for:} Upward trends, spikes, seasonal patterns
    \item \textbf{Action:} If spike, investigate bonuses or back pay
\end{itemize}

\textbf{Plot 2: Cost Forecast}
\begin{itemize}
    \item \textbf{Look for:} Predicted trend (up/down/stable)
    \item \textbf{Action:} Use for budget planning next 6 months
    \item \textbf{Note:} Confidence interval widens over time (uncertainty)
\end{itemize}

\textbf{Plot 3: Cost Per Employee}
\begin{itemize}
    \item \textbf{Look for:} Increasing = hiring expensive staff, Decreasing = cost reduction
    \item \textbf{Action:} Compare to industry benchmarks
\end{itemize}

\textbf{Plot 4: Cost Breakdown}
\begin{itemize}
    \item \textbf{Look for:} Proportion of net vs taxes
    \item \textbf{Action:} If taxes increase, review tax optimization
\end{itemize}

\textbf{Plot 5: Cost Distribution}
\begin{itemize}
    \item \textbf{Look for:} Outliers (dots beyond whiskers), skewness
    \item \textbf{Action:} Outliers = investigate salary equity
\end{itemize}

\textbf{Plot 6: Top Employees}
\begin{itemize}
    \item \textbf{Look for:} Top 10 most expensive employees
    \item \textbf{Action:} Retention planning for high-cost employees
\end{itemize}

\textbf{Plot 7: Gross vs Net}
\begin{itemize}
    \item \textbf{Look for:} Gap between gross and net (tax burden)
    \item \textbf{Action:} Individual salary review, tax optimization
\end{itemize}

\textbf{Plot 8: Headcount Trend}
\begin{itemize}
    \item \textbf{Look for:} Hiring/firing patterns, growth rate
    \item \textbf{Action:} Correlate with revenue to assess efficiency
\end{itemize}

\textbf{Plot 9: Burn Rate Dashboard}
\begin{itemize}
    \item \textbf{Look for:} Monthly variability, cumulative total, runway
    \item \textbf{Action:} If runway <6 months, urgent cost reduction needed
\end{itemize}

\textbf{Plot 10: Year-over-Year}
\begin{itemize}
    \item \textbf{Look for:} Seasonal patterns, growth rate
    \item \textbf{Action:} Compare same month across years for trend analysis
\end{itemize}

\section{Workflow 3: Modify Existing Mappings}

\subsection{Step-by-Step Process}

\textbf{Step 1: Open Source Workbook}

\begin{enumerate}
    \item Open the payroll workbook (e.g., \texttt{Pay\_emploier\_sept25.xlsm})
    \item Navigate to \textbf{CustomMap} sheet (created by first mapping)
\end{enumerate}

\textbf{Step 2: Review Mappings}

\begin{table}[H]
\centering
\footnotesize
\begin{tabular}{|l|l|l|}
\hline
\textbf{Destination\_Label} & \textbf{Mapping\_Type} & \textbf{KeysOps} \\
\hline
salaire\_brut & formula & +::SALAIRE\_BASE::label \\
salaire\_net & formula & +::BRUT::label;;-::COT::label \\
prime\_total & formula & +::PRIME1::label;;+::PRIME2::label \\
\hline
\end{tabular}
\caption{CustomMap Sheet Structure}
\end{table}

\textbf{Step 3: Edit Formula}

\begin{enumerate}
    \item Find row with destination you want to change
    \item Edit KeysOps column:
    \begin{itemize}
        \item Format: \texttt{operator::key::type;;operator::key::type}
        \item Example: \texttt{+::BASE::label;;+::PRIME::label}
    \end{itemize}
    \item Save workbook
\end{enumerate}

\textbf{Step 4: Re-run Mapping}

\begin{enumerate}
    \item Open \texttt{Control\_Panel\_New.xlsm}
    \item Click \textbf{Button 1: Import \& Map Data}
    \item GUI loads existing mappings from CustomMap sheet
    \item Click \textbf{"Execute Mapping"} (only modified columns will recalculate)
\end{enumerate}

\section{Workflow 4: Generate Synthetic Data}

\subsection{Purpose}

When real payroll data is too small or sensitive, generate synthetic data for:
\begin{itemize}
    \item Model training
    \item Algorithm testing
    \item Public demonstrations
    \item Academic papers
\end{itemize}

\subsection{Step-by-Step Process}

\textbf{Step 1: Configure Generator}

Edit \texttt{src/data generation/generate\_synthetic\_payroll.py}:

\begin{lstlisting}[style=pythonstyle]
# Configuration
num_employees = 50        # Number of synthetic employees
num_months = 24          # Number of months to generate
base_salary_mean = 15000 # Average base salary
base_salary_std = 3000   # Salary variation

# Growth parameters
monthly_growth_rate = 0.02  # 2% monthly growth
turnover_rate = 0.05        # 5% monthly turnover
\end{lstlisting}

\textbf{Step 2: Run Generator}

\begin{lstlisting}
python src/data generation/generate_synthetic_payroll.py
\end{lstlisting}

\textbf{Step 3: Output}

\begin{itemize}
    \item Saves to \texttt{outputs/synthetic\_payroll\_long.csv}
    \item Same structure as real data (long format)
    \item Can be used directly with dashboard and forecasting
\end{itemize}

\subsection{Synthetic Data Features}

\begin{enumerate}
    \item \textbf{Realistic distributions:} Normal distribution for salaries
    \item \textbf{Temporal trends:} Growth over time
    \item \textbf{Employee churn:} Hiring/terminations
    \item \textbf{Noise:} Random monthly variations
    \item \textbf{Correlations:} Gross/net relationship preserved
\end{enumerate}

\section{Workflow 5: Export for Analysis}

\subsection{Export Long Format CSV}

\textbf{Purpose:} Create CSV file for analysis in Python, R, or other tools.

\textbf{Steps:}

\begin{enumerate}
    \item Ensure \texttt{data\_sorted.xlsx} exists (run Workflow 1 if needed)
    \item Run conversion script:
    \begin{lstlisting}
python scripts/run_load_data.py
    \end{lstlisting}
    
    \item Output: \texttt{outputs/payroll\_long.csv}
    
    \item Verify structure:
    \begin{lstlisting}[style=pythonstyle]
import pandas as pd
df = pd.read_csv('outputs/payroll_long.csv')
print(df.head())
print(df.columns)
    \end{lstlisting}
\end{enumerate}

\textbf{Expected Columns:}
\begin{itemize}
    \item \texttt{month} - Date (YYYY-MM-DD)
    \item \texttt{employee\_id} - Unique identifier
    \item \texttt{salaire\_brut} - Gross salary
    \item \texttt{salaire\_net} - Net salary
    \item Additional columns as defined in mappings
\end{itemize}

\subsection{Export Forecast Results}

\begin{enumerate}
    \item Run forecasting script:
    \begin{lstlisting}
python scripts/run_forecast.py
    \end{lstlisting}
    
    \item Output: \texttt{outputs/forecast\_ar2.csv}
    
    \item Structure:
    \begin{table}[H]
    \centering
    \begin{tabular}{cc}
    \toprule
    \textbf{forecast\_month} & \textbf{predicted\_cost} \\
    \midrule
    1 & 145000 \\
    2 & 148000 \\
    3 & 150000 \\
    ... & ... \\
    \bottomrule
    \end{tabular}
    \caption{Forecast CSV Structure}
    \end{table}
\end{enumerate}

\section{Best Practices}

\subsection{Data Quality}

\begin{enumerate}
    \item \textbf{Consistent naming:} Use same labels across months
    \item \textbf{No missing IDs:} Ensure Row 7 has IDs for all employees
    \item \textbf{Numeric values:} Avoid text in data cells
    \item \textbf{Date format:} Use consistent month sheet names
\end{enumerate}

\subsection{Mapping Formulas}

\begin{enumerate}
    \item \textbf{Test incrementally:} Define one formula at a time, verify
    \item \textbf{Use codes when possible:} More reliable than labels
    \item \textbf{Document mappings:} Add comments in CustomMap sheet
    \item \textbf{Version control:} Save source workbook after each mapping session
\end{enumerate}

\subsection{Visualization}

\begin{enumerate}
    \item \textbf{Regular monitoring:} Run dashboard weekly/monthly
    \item \textbf{Share insights:} Export plots (right-click plot > Save As)
    \item \textbf{Compare periods:} Use YoY plot to track progress
    \item \textbf{Act on anomalies:} Investigate spikes/drops immediately
\end{enumerate}

\subsection{Performance}

\begin{enumerate}
    \item \textbf{Close other Excel files:} Reduces xlwings overhead
    \item \textbf{Use "Modified Only":} When re-running mappings
    \item \textbf{Batch processing:} Process multiple months at once
    \item \textbf{Cache results:} Save intermediate CSVs to avoid recomputation
\end{enumerate}

\section{Command Line Reference}

\subsection{Manual Execution}

If buttons don't work, run scripts manually:

\begin{lstlisting}
# Import and mapping
python scripts/run_excel.py

# Dashboard
python scripts/run_dashboard.py

# Convert to long format
python scripts/run_load_data.py

# Forecasting
python scripts/run_forecast.py

# Plotting (individual plots)
python scripts/run_plot.py
\end{lstlisting}

\subsection{Python Environment}

Check installed packages:

\begin{lstlisting}
pip list | grep pandas
pip list | grep xlwings
pip list | grep statsmodels
pip list | grep matplotlib
\end{lstlisting}

Update packages:

\begin{lstlisting}
pip install --upgrade pandas xlwings statsmodels matplotlib
\end{lstlisting}

\section{Frequently Asked Questions}

\textbf{Q: Can I use different month sheet names?}

A: Yes, but update \texttt{src/core/config.py}:
\begin{lstlisting}[style=pythonstyle]
MONTH_SHEETS = ["Jan", "Feb", "Mar", ...]
\end{lstlisting}

\textbf{Q: How do I add a new destination column?}

A: Edit target workbook header row, then define formula in GUI.

\textbf{Q: Can I use division in formulas?}

A: Yes, use \texttt{/} operator. Example: \texttt{+TOTAL label /HEURES label}

\textbf{Q: What if I have multiple years of data?}

A: Name sheets like "2024-01", "2024-02", etc. Update config accordingly.

\textbf{Q: How do I delete a mapping?}

A: Open CustomMap sheet, delete the row, save workbook.

\textbf{Q: Can I run this on Mac/Linux?}

A: Python scripts yes, Excel macros no (VBA is Windows-only). Use LibreOffice Calc with Python scripts directly.

\textbf{Q: How do I backup my mappings?}

A: CustomMap sheet is stored in source workbook. Back up the entire workbook file.

\textbf{Q: What's the maximum number of employees?}

A: No hard limit, but performance degrades above 1000 employees. Consider splitting data.

\textbf{Q: Can I forecast more than 6 months?}

A: Yes, change \texttt{forecast\_steps} parameter, but accuracy decreases with longer horizons.

\textbf{Q: How do I interpret confidence intervals?}

A: Wider interval = more uncertainty. If interval is too wide (>50\% of forecast), don't rely on forecast.
